\subsection{同核复用和跨核传播的集合论表示}
场景 B 中“同核子图同一 Task”改变了切图的边际成本。对内部张量 $x$，以生产核心 $c_p(x)$ 和去重后的消费核心集合 $\mathcal C(x;p)$ 描述传播需求。若 $c_p(x)\in\mathcal C(x;p)$，在生产核心上继续消费无需因子图边界而复制到 DDR；其它消费核心分别需要获得该张量。一个仅用于静态排序的跨核读取机会数是
\begin{equation}\label{eq:q2-read-opportunity}
R_x^{\rm remote}(p)=
\sum_{c=1}^{k}\mathbf 1\{c\in\mathcal C(x;p),\ c\ne c_p(x)\}
=n_{\rm remote}(x;p).
\end{equation}
若把 $x$ 的三个消费者映射到同一个远端核心，$R_x^{\rm remote}=1$；把它们映射到三个不同的远端核心则为 $3$。从式\eqref{eq:q2-read-opportunity}可推出，按算子边数惩罚跨核访问会高估第一个布局的读取需求，并可能错误偏好把互不相关计算都聚在源核。对外部输入则没有 $c_p(x)$，应从不同消费核心的数量出发，而不能凭空添加一次“源核心写出”。

然而 $b_x R_x^{\rm remote}$ 仍\emph{不是}精确额外搬运：一次跨核传输可能要求源端写出，原始图可能已含相关 COPY，多个消费者有不同的私有缓存生存区间，核内启发式也可能额外换入。严谨比较应先在两方案 $p,p'$ 上分别构造官方 COPY 事件集合 $\mathcal E_D(p),\mathcal E_D(p')$，计算 $D^B(p),D^B(p')$，再将其作为次指标。静态 $R_x^{\rm remote}$ 的职责只是决定哪些共享张量的消费者有合并到同核的潜力。本文把“跨核目标集合减少”与“正式额外搬运必减少”明确分开，避免把代理变量写成赛题评分变量。

如果一个远端核心内的两个消费算子之间插入很多与 $x$ 无关的工作，$x$ 的逻辑活跃区间变长。设首次使用在 $t_1$，最后一次使用在 $t_2$，那么理想完全驻留占用面积为 $b_x(t_2-t_1)$；把两次使用靠近可能缩小这部分面积，并在某些拥塞情况下减少溢出。但是移动顺序还会改变其它张量的生命周期，不能仅按一个张量的面积贪心。对候选顺序可定义加权峰值代理 $\max_tL^{\rm ideal}_{c,r}(t)$，并与跨核读取机会数一起观察：如果为了减少 $R_x^{\rm remote}$ 而令峰值显著超过 L1/UB 容量，便应保留原布局参加官方对照，而不是无条件接受聚集。

\subsection{同步释放的精确落点与链式传播}
场景 B 的 $500$ 周期是源端特定 	exttt{COPY\_OUT} 完成与目标端对应 	exttt{COPY\_IN} 可发射之间的约束。若源任务还有后续的无关计算，目标不必等源 Task 完成；若目标核心在释放期间执行另一条独立分支，释放不一定增加 Makespan。设某目标读取的源写出在 $f_{
m out}$ 完成，目标 DDR Pipe 最早在 $a_{
m pipe}$ 空闲，其数据前驱允许的最早时刻为 $a_{
m pred}$，则它的发射下界为
\begin{equation}\label{eq:q2-copy-start}
s_{
m in}\ge\max\{f_{
m out}+500,\ a_{
m pipe},\ a_{
m pred}\}.
\end{equation}
写成 $f_{
m out}+500+a_{
m pipe}+a_{
m pred}$ 显然会把可并行等待相加。若目标之后有矩阵操作 $v$ 必须等该读取完成，其开始还要满足 $s_v\ge f_{
m in}$，但同核心的其它 Vector 操作只服从自己的依赖与 Pipe 顺序。式\eqref{eq:q2-copy-start}与事件释放式\eqref{eq:event-release}共同决定同步开销如何传播到最终输出。这一传播链较“每个跨核边附加常数 500”更准确。

在相同图中，单纯迁移一个子图可能改变多条同步边的方向与数量。设子图 $S$ 原在核心 $a$，迁至核心 $b$；其前驱生产核心集合为 $P(S)$，后继消费核心集合为 $Q(S)$。若 $P(S)$ 中许多生产者在 $b$，迁移可删去若干输入同步；但如果 $Q(S)$ 中多数消费者仍在 $a$，输出同步可能增加。潜在净变化可以先用集合差
\begin{equation}\label{eq:q2-sync-proxy}
\Delta\widehat N_{\times}(S:a\to b)=
\sum_{x:\,C(x)\cap S\ne\varnothing}\!
\bigl(R_x^{\rm remote}(p')-R_x^{\rm remote}(p)\bigr)
\end{equation}
估计，而非把所有边计一次。但即使 $\Delta\widehat N_{\times}<0$，也可能只删除非关键分支上的同步，或者增加更长的同核内存活跃区间。真实接受规则仍取完整官方 $T^B$ 的差值。这里的集合表达是候选筛选的合理数学依据，未假定每少一次同步就能赚到整整 $500$ 周期。

\subsection{划分与映射联合搜索的可解释性}
把边界、核心和局部顺序分成完全独立的三次优化并不充分。给定一个边界动作 $a_\pi$，其潜在收益可能只有在随后迁核 $a_m$ 并插入到释放空隙 $a_\sigma$ 后才能显现；只评价中途的 $a_\pi$，会误把整个联合动作淘汰。为避免对三维组合穷举，本文只构造少量由共同瓶颈触发的联动候选。令 $\mathcal R$ 为关键输出附近且共享某些张量的有限区域，候选写成
\begin{equation}\label{eq:q2-joint-transform}
p'=\operatorname{RepairOrder}
\bigl(\operatorname{MoveCore}(\operatorname{ShiftBoundary}(p,\mathcal R),\mathcal R)\bigr).
\end{equation}
每一步允许“保持原值”以覆盖单动作，但总提案数和官方机会数设硬上限。候选必须在完整动作结束后统一判定是否合法，不能用某个中间状态的差周期代表最终联合动作；同时必须保留原方案，避免错误代理将有效布局丢失。实际 r07 交付仍采用式\eqref{eq:q2-neighborhood}的有限普通联合迁移与局部顺序修复，式\eqref{eq:q2-joint-transform}描述更一般的可解释邻域形状，不宣称交付版遍历了其中所有边界组合。

有限联合候选的评价顺序应兼顾结构差异。例如两个方案虽然算子分组不同，但如果由官方展开后得到完全相同的 Task、COPY 及核上顺序，其官方结果也会一致，重复占用评价机会没有信息价值。规范指纹首先包含原图哈希、核数和场景参数，然后包含每个可切算子的子图标识、每核有序子图列；仅凭生成过程名称或代理分数无法识别同一方案。若不同标识只是子图编号的置换，应在指纹前规范化编号。这个去重过程不改变算法的数学可行集，却直接提高固定预算中真正不同候选的覆盖率；它是“更多算子”之外的一种求解效率改进。

\subsection{场景转换的独立对照原则}
问题二相对问题一增添同核复用并改变同步语义，因此 $T^B_{i,k}<T^A_{i,k}$ 不能全部归因于“问题二算法更聪明”。若要分离场景机制和搜索机制，应对\emph{同一个} $p$ 同时计算 $T^A(p)$ 与 $T^B(p)$，再在 B 中比较移植 $p^A$ 与新选 $p^B$。形式上，若符号与合法性允许，可写总差
\begin{equation}\label{eq:q2-scene-alg-decomp}
T^A(p^A)-T^B(p^B)=
\bigl[T^A(p^A)-T^B(p^A)\bigr]
+\bigl[T^B(p^A)-T^B(p^B)\bigr].
\end{equation}
第一项是固定 A 方案的场景变化，第二项是 B 场景中的重新选方案。若某移植方案不满足 B 的物理约束，应先执行正式转换与合法化，不能强行套式。论文中的三问主表主要报告各场景最终效果，不把式\eqref{eq:q2-scene-alg-decomp}两项写成已经全量测得的独立因果估计；这个分解为将来做更完整的统一消融提供准确口径。它也解释问题二为何必须重新运行官方评价，而不能只把 A 的 $1000$ 改成 $500$ 再报告预计收益。
